%% kapa: exam=2 round=26 year=2015 period=1 subject=appraisal_practice subject_ko=감정평가실무 seq=1 number=1 points=40
%% source: Q-Net 감정평가사 기출문제 (cst00309) artlSeq=5149601 file=extracted/1008/5149601/2015년 제26회 감정평가사 2차시험 기출문제/1교시(감정평가실무).pdf
%% conversion: kordoc 4.24.51 -> tools/gen2.py (auto); review=proofread-gpt-sol; flags=none
\begin{longquestion}{【문제 1】}
한국○○공사는 보유중인 부동산을 매각하기 위해 김공정 감정평가사에게 일반거래(시가참고) 목적의 감정평가를 의뢰하였다. 관련법규 및 이론을 참작하고 제시된 자료를 활용하여 다음의 물음에 답하시오. (40점) \par
\medskip\noindent 물음 1) 본 감정평가에 적용할 층별효용지수를 산정하시오. (10점) \par
\medskip\noindent 물음 2) 시산가액 조정을 통해 감정평가액을 구하시오. (30점) \par
<자료 1> 기본적 사항 \par
1. 감정평가 의뢰내역 \par
{\small\centering
\resizebox{\ifdim\width>\linewidth\linewidth\else\width\fi}{!}{%
\begin{tabular}{|c|c|c|c|c|c|c|}
\hline
\makecell{기호} & \makecell{소재지 지번} & \makecell{층} & \makecell{호수} & \makecell{전용면적\\(m\textsuperscript{2})} & \makecell{공용면적\\(m\textsuperscript{2})} & \makecell{전체면적\\(m\textsuperscript{2})} \\
\hline
\makecell{1} & \makecell{서울시 A구 B동 00번지} & \makecell{1층} & \makecell{101} & \makecell{1,350} & \makecell{1,650} & \makecell{3,000} \\
\hline
\makecell{2} & \makecell{〃} & \makecell{2층} & \makecell{201} & \makecell{1,215} & \makecell{1,485} & \makecell{2,700} \\
\hline
\makecell{3} & \makecell{〃} & \makecell{3층} & \makecell{301} & \makecell{1,215} & \makecell{1,485} & \makecell{2,700} \\
\hline
\makecell{4} & \makecell{〃} & \makecell{4층} & \makecell{401} & \makecell{1,100} & \makecell{1,600} & \makecell{2,700} \\
\hline
\makecell{5} & \makecell{〃} & \makecell{5층} & \makecell{501} & \makecell{1,215} & \makecell{1,485} & \makecell{2,700} \\
\hline
\makecell{6} & \makecell{〃} & \makecell{6층} & \makecell{601} & \makecell{900} & \makecell{1,100} & \makecell{2,000} \\
\hline
\makecell{7} & \makecell{〃} & \makecell{지1층} & \makecell{B101} & \makecell{2,250} & \makecell{2,750} & \makecell{5,000} \\
\hline
\end{tabular}}\par}

2. 기준시점: 2015. 08. 20 \par
3. 기준가치: 시장가치 \par
4. 평가목적: 일반거래(시가참고) \par
<자료 2> 지역분석 자료 \par
1. 본건이 위치하고 있는 지역은 새롭게 조성된 상업 및 업무지대로 토지 및 업무시설(집합건물)의 평가사례 및 거래사례가 풍부함 \par
2. 기준시점 현재 해당지역의 업무시설가격은 2013년 1분기 대비 소폭 상승하였으나 해당기간동안 업무시설가격은 상승과 하락을 반복하였음 \par
3. 사무실 또는 상가는 층별 각기 다른 가격격차를 보이고 있는데, 고객의 이용에 따른 편의성, 접근성, 수익성 등에 따른 것으로 판단됨. 한편, 지하철역과의 거리에 따른 가격격차도 확인할 수 있었음. 그리고 지하철역과의 거리에 따른 가격격차는 업무용 토지가격에서도 확인할 수 있었음 \par
4. 상기의 지역분석은 탐문조사, 평가사례 및 거래사례 등을 이용하여 분석한 것으로 보다 상세한 해당지역의 가치형성요인을 분석하기 위해 기준시점으로부터 6개월 이내 자료를 이용하여 계량분석을 실시함 \par
<자료 3> 가치형성요인의 계량분석 \par
1. 헤도닉가격모형을 이용하여 해당지역의 가치형성요인을 분석함 \par
2. 토지만의 거래사례를 이용하여 업무용 토지가격을 종속변수로 한 모형을 추정한 결과 모형의 설명력은 0.875(수정된 R제곱)이고, F-value는 629.430으로 나타남. 다음은 분석내용임 \par
1) 설명변수로 채택된 업무용토지의 면적은 5 \% 유의수준에서 5,000 m\textsuperscript{2}∼ 16,500 m\textsuperscript{2} 면적의 토지는 5,000 m\textsuperscript{2}미만 면적의 토지에 대해 다른 조건이 일정할 때 가격측면에서 약 10 \% 우세한 것으로 나타났으며 이는 통계적으로 유의함 \par
2) 설명변수로 채택된 본건과 지하철역과의 거리는 5 \% 유의수준에서 유의 했는데, 지하철과의 거리가 0.5 km∼1 km인 토지는 0.5 km미만의 토지에 대해 다른 조건이 일정할 때 가격측면에서 3 \% 열세한 것으로 나타났으나 1 km 초과 토지는 통계적으로 유의하지 않음 \par
3. 집합건물인 업무시설(사무실)의 거래사례를 이용하여 업무시설가격을 종속변수로 한 모형을 추정한 결과 모형의 설명력은 0.825(수정된 R제곱)이고, F-value는 523.257로 나타남. 다음은 분석내용임 \par
1) 설명변수로 채택된 층의 경우 본 계량모형상의 1층 가격은 약 3,500,000 원/m\textsuperscript{2} 정도임. 한편, 다른 조건이 모두 동일한 경우 지하1층은 1 \% 유의 수준에서 880,000 원/m\textsuperscript{2} 정도가 1층에 비해 가격이 낮게 나타남. 한편, 2층부터 6층까지는 5 \% 유의수준에서 2층은 410,000 원/m\textsuperscript{2}, 3층은 295,000 원/m\textsuperscript{2}, 4층은 385,000 원/m\textsuperscript{2}, 5층은 350,000 원/m\textsuperscript{2}, 6층 이상은 400,000 원/m\textsuperscript{2} 정도가 1층에 비해 가격이 낮게 나타났으며 이는 통계적으로 유의함 \par
2) 설명변수로 채택된 지하철역까지의 거리 변수는 5 \% 유의수준에서 유의했고 지하철역에서 멀어질수록 업무시설가격은 하락(-)함 \par
3) 설명변수로 채택된 전용율의 경우 전체면적이 통제된 상태에서 전용면적의 증가는 업무시설가격에 긍정적인(+) 효과를 미쳤고 1 \% 유의수준에서 유의한 것으로 나타남. 구체적으로 전용율 45 \% 미만의 업무시설은 전용율 45 \% 이상의 업무시설에 비해 가격측면에서 약 3 \% 열세한 것으로 나타남. 한편, 전용면적이 통제된 상태에서 공용면적의 증가는 통계적으로 유의하지 않음 \par
4) 설명변수로 채택된 업무시설의 전체면적은 1 \% 유의수준에서 4,000 m\textsuperscript{2} ∼8,000 m\textsuperscript{2} 면적의 업무시설은 4,000 m\textsuperscript{2} 미만 면적의 업무시설에 비해 다른 조건이 일정할 때 가격측면에서 약 5 \% 열세한 것으로 나타났으며 통계적으로 유의함 \par
<자료 4> 개별요인분석 자료 \par
1. 대상부동산은 건물전체가 의뢰된 건으로 총 7개의 집합건물로 구성되어 있고, 적법한 절차를 걸쳐 완공된 상태임 \par
2. 토지 \par
{\small\centering
\resizebox{\ifdim\width>\linewidth\linewidth\else\width\fi}{!}{%
\begin{tabular}{|c|c|}
\hline
\makecell{구분} & \makecell{내 용} \\
\hline
\makecell{지목 및 면적} & \makecell{대 / 3,637 m\textsuperscript{2}} \\
\hline
\makecell{위치 및\\주위환경} & \makecell{서울시 A구 B동 소재 H백화점 북측인근에 소재하며\\인근은 상업 및 업무지대임} \\
\hline
\makecell{도로 및\\교통환경} & \makecell{본건 남측으로 지하철 K역(본건과의 거리: 0.62 km)이 위치해\\있고 동측 및 남측인근에 일반버스정류장이 소재하는\\등 교통사항은 양호함} \\
\hline
\makecell{지형, 지세\\및 이용 상황} & \makecell{남측 및 서측은 인접도로와 등고평탄한 정방형 토지로\\기준시점 현재 업무시설부지로 이용중임} \\
\hline
\makecell{토지이용\\계획사항} & \makecell{도시지역, 준주거지역, 제1종 지구단위계획구역, 대로2류\\(폭 30 M∼35 M)(접함), 중로2류(폭 15 M∼20 M)(접함)} \\
\hline
\end{tabular}}\par}

3. 건물 \par
{\small\centering
\resizebox{\ifdim\width>\linewidth\linewidth\else\width\fi}{!}{%
\begin{tabular}{|c|c|}
\hline
\makecell{구분} & \makecell{내 용} \\
\hline
\makecell{구조 및\\용도} & \makecell{철골철근콘크리트조, 슬래브지붕 / 업무시설} \\
\hline
\makecell{건축면적\\및 연면적} & \makecell{2,228 m\textsuperscript{2} / 20,800 m\textsuperscript{2}} \\
\hline
\makecell{층수 등} & \makecell{ㆍ지상6층(지상1층～지상6층: 업무시설)\\ㆍ지하3층(지하1층 : 업무시설, 지하2, 3층 : 주차장, 기계실, 전기실 등)} \\
\hline
\makecell{사용승인일} & \makecell{2013년 7월 20일} \\
\hline
\makecell{주차 및\\부대설비} & \makecell{140대(주차), 엘리베이터 4대,\\에스컬레이터 2대, 소화설비, 냉난방설비 등} \\
\hline
\end{tabular}}\par}

<자료 5> 집합건물 거래사례(소재지: 서울시 A구 B동) \par
{\small\centering
\resizebox{\ifdim\width>\linewidth\linewidth\else\width\fi}{!}{%
\begin{tabular}{|c|c|c|c|c|c|c|c|c|}
\hline
\makecell{기호} & \makecell{지번} & \makecell{층} & \makecell{호수} & \makecell{전용면적\\(m\textsuperscript{2})} & \makecell{전체\\면적\\(m\textsuperscript{2})} & \makecell{거래가격\\(원)} & \makecell{용도} & \makecell{거래시점} \\
\hline
\makecell{A} & \makecell{625} & \makecell{3} & \makecell{301} & \makecell{1,350} & \makecell{3,000} & \makecell{8,500,000,000} & \makecell{사무소} & \makecell{2015.01.10} \\
\hline
\makecell{B} & \makecell{670} & \makecell{4} & \makecell{401} & \makecell{1,125} & \makecell{2,500} & \makecell{6,780,000,000} & \makecell{사무소} & \makecell{2015.02.10} \\
\hline
\makecell{C} & \makecell{710} & \makecell{5} & \makecell{501} & \makecell{1,080} & \makecell{2,700} & \makecell{6,350,000,000} & \makecell{연구소} & \makecell{2015.03.10} \\
\hline
\makecell{D} & \makecell{720} & \makecell{지1} & \makecell{B101} & \makecell{2,080} & \makecell{5,200} & \makecell{11,000,000,000} & \makecell{사무소} & \makecell{2015.04.10} \\
\hline
\end{tabular}}\par}

1. 지하철역까지의 거리 \par
ㆍ기호A: 0.55 km, 기호B: 0.63 km, 기호C: 0.60 km, 기호D: 1.80 km \par
2. 용도지역 \par
ㆍ기호A: 일반상업, 기호B: 준주거, 기호C: 준주거, 기호D: 준주거 \par
3. 가치형성요인의 경우 상기자료 외에는 대상과 사례가 동일함 \par
<자료 6> 층별효용지수 참고자료 \par
1. 계량분석외 본건에 인접한 유사 건물의 평가사례 및 실무기준해설서상의 층별효용지수 관련 자료는 다음과 같음 \par
{\small\centering
\resizebox{\ifdim\width>\linewidth\linewidth\else\width\fi}{!}{%
\begin{tabular}{|c|c|c|c|c|c|}
\hline
\multirow{2}{*}{\makecell{층}} & \multirow{2}{*}{\makecell{평가사례1\\(2013.03.27)}} & \multirow{2}{*}{\makecell{평가사례2\\(2014.09.15)}} & \multirow{2}{*}{\makecell{평가사례3\\(2015.07.20)}} & \multicolumn{2}{|c|}{\makecell{실무기준해설서}} \\
\cline{5-6}
 &  &  &  & \makecell{A형} & \makecell{B형} \\
\hline
\makecell{지상5층이상} & \makecell{85} & \makecell{87} & \makecell{90} & \makecell{42} & \makecell{51} \\
\hline
\makecell{지상4층} & \makecell{92} & \makecell{87} & \makecell{90} & \makecell{45} & \makecell{51} \\
\hline
\makecell{지상3층} & \makecell{83} & \makecell{87} & \makecell{90} & \makecell{50} & \makecell{51} \\
\hline
\makecell{지상2층} & \makecell{85} & \makecell{87} & \makecell{90} & \makecell{60} & \makecell{51} \\
\hline
\makecell{지상1층} & \makecell{100} & \makecell{100} & \makecell{100} & \makecell{100} & \makecell{100} \\
\hline
\makecell{지하1층} & \makecell{-} & \makecell{-} & \makecell{-} & \makecell{44} & \makecell{44} \\
\hline
\end{tabular}}\par}

<자료 7> 공시지가 표준지 \par
1. 공시기준일: 2015년 1월 1일 \par
2. 인근지역내 공시지가 표준지 내역(소재지: 서울시 A구 B동) \par
{\small\centering
\resizebox{\ifdim\width>\linewidth\linewidth\else\width\fi}{!}{%
\begin{tabular}{|c|c|c|c|c|c|c|c|c|c|}
\hline
\makecell{기호} & \makecell{지번} & \makecell{지목} & \makecell{면적\\(m\textsuperscript{2})} & \makecell{이용\\상황} & \makecell{용도지역} & \makecell{도로\\교통} & \makecell{형상\\지세} & \makecell{공시지가\\(원/m\textsuperscript{2})} & \makecell{지하철역과의\\거리} \\
\hline
\makecell{가} & \makecell{630} & \makecell{대} & \makecell{5,179.2} & \makecell{업무용} & \makecell{일반상업} & \makecell{중로\\각지} & \makecell{세장형\\평지} & \makecell{3,000,000} & \makecell{0.5 km} \\
\hline
\makecell{나} & \makecell{681} & \makecell{대} & \makecell{3,329.2} & \makecell{업무용} & \makecell{준주거} & \makecell{광대\\소각} & \makecell{정방형\\평지} & \makecell{3,250,000} & \makecell{1.5 km} \\
\hline
\makecell{다} & \makecell{702} & \makecell{대} & \makecell{9,038.0} & \makecell{업무용} & \makecell{준주거} & \makecell{중로\\각지} & \makecell{정방형\\평지} & \makecell{2,800,000} & \makecell{0.7 km} \\
\hline
\end{tabular}}\par}

<자료 8> 토지거래사례 \par
(소재지: 서울시 A구 B동, 거래시점: 2015년 7월 10일) \par
{\small\centering
\resizebox{\ifdim\width>\linewidth\linewidth\else\width\fi}{!}{%
\begin{tabular}{|c|c|c|c|c|c|c|c|c|}
\hline
\makecell{기호} & \makecell{지번} & \makecell{지목} & \makecell{면적\\(m\textsuperscript{2})} & \makecell{이용\\상황} & \makecell{용도\\지역} & \makecell{도로\\교통} & \makecell{형상\\지세} & \makecell{거래단가\\(원/m\textsuperscript{2})} \\
\hline
\makecell{1} & \makecell{698} & \makecell{대} & \makecell{3,500} & \makecell{업무용} & \makecell{준주거} & \makecell{광대\\한면} & \makecell{사다리형\\평지} & \makecell{5,800,000} \\
\hline
\end{tabular}}\par}

※ 지하철역과의 거리 : 0.3 km \par
<자료 9> 재조달원가 자료(기준시점 기준) \par
{\small\centering
\resizebox{\ifdim\width>\linewidth\linewidth\else\width\fi}{!}{%
\begin{tabular}{|c|c|c|c|c|c|}
\hline
\makecell{기호} & \makecell{용도} & \makecell{구조} & \makecell{급수} & \makecell{단가\\(원/m\textsuperscript{2})} & \makecell{내용 년수} \\
\hline
\makecell{1} & \makecell{사무실} & \makecell{철골철근콘크리트조\\슬래브지붕 (6층이하)} & \makecell{1급} & \makecell{1,540,000} & \makecell{55 (50-60)} \\
\hline
\makecell{2} & \makecell{사무실} & \makecell{철근콘크리트조\\슬래브지붕 (6층이하)} & \makecell{1급} & \makecell{1,400,000} & \makecell{55 (50-60)} \\
\hline
\end{tabular}}\par}

<자료 10> 수익관련자료 \par
1. 임대사례[소재지: 서울시 A구 B동, 임대면적: 3,000 m\textsuperscript{2}(전체)] \par
{\small\centering
\resizebox{\ifdim\width>\linewidth\linewidth\else\width\fi}{!}{%
\begin{tabular}{|c|c|c|c|c|c|c|c|c|}
\hline
\makecell{기호} & \makecell{지번} & \makecell{층} & \makecell{호수} & \makecell{전용율} & \makecell{보증금\\(원/m\textsuperscript{2})} & \makecell{월임대료\\(원/m\textsuperscript{2})} & \makecell{월관리비\\(원/m\textsuperscript{2})} & \makecell{지하철역과의 거리} \\
\hline
\makecell{1} & \makecell{699} & \makecell{1} & \makecell{101} & \makecell{45 \%} & \makecell{100,000} & \makecell{10,000} & \makecell{6,000} & \makecell{0.63 km} \\
\hline
\end{tabular}}\par}

2. 임대사례는 기준시점 현재 당해 업무지대의 평균적인 사례임 \par
3. 공실 및 대손충당금 : 가능조소득의 10 \% \par
4. 영업경비: 연간관리비의 80 \% \par
5. 보증금, 월임대료, 월관리비는 전체면적기준임 \par
6. 종합환원이율 : 연 4.0 \%, 보증금운용이율 : 연 3.0 \% \par
7. 가치형성요인의 경우 상기자료 외에는 대상과 사례가 동일함 \par
<자료 11> 기타자료 \par
1. 도로접면 격차율 \par
{\small\centering
\resizebox{\ifdim\width>\linewidth\linewidth\else\width\fi}{!}{%
\begin{tabular}{|c|c|c|c|c|c|}
\hline
\makecell{구분} & \makecell{광대한면} & \makecell{중로한면} & \makecell{소로한면} & \makecell{세로(가)} & \makecell{세로(불)} \\
\hline
\makecell{격차율} & \makecell{1.00} & \makecell{0.90} & \makecell{0.88} & \makecell{0.85} & \makecell{0.70} \\
\hline
\end{tabular}}\par}

※ 도로접면이 각지인 경우는 한면에 접하는 경우에 비해 1 \% 우세함 \par
2. 형상 격차율 \par
{\small\centering
\resizebox{\ifdim\width>\linewidth\linewidth\else\width\fi}{!}{%
\begin{tabular}{|c|c|c|c|c|c|}
\hline
\makecell{구분} & \makecell{정방형} & \makecell{가장형} & \makecell{세장형} & \makecell{사다리형} & \makecell{부정형} \\
\hline
\makecell{격차율} & \makecell{1.00} & \makecell{0.99} & \makecell{0.99} & \makecell{0.98} & \makecell{0.90} \\
\hline
\end{tabular}}\par}

3. 시점수정 자료 \par
1) 지가변동률(단위: \%) \par
{\small\centering
\resizebox{\ifdim\width>\linewidth\linewidth\else\width\fi}{!}{%
\begin{tabular}{|c|c|c|}
\hline
\multirow{2}{*}{\makecell{기간}} & \multicolumn{2}{|c|}{\makecell{용도지역}} \\
\cline{2-3}
 & \makecell{주거지역} & \makecell{상업지역} \\
\hline
\makecell{2015. 01. 01 ∼ 2015. 07. 31(누계)} & \makecell{0.435} & \makecell{0.420} \\
\hline
\makecell{2015. 07. 01 ∼ 2015. 07. 31(당월)} & \makecell{0.020} & \makecell{0.010} \\
\hline
\end{tabular}}\par}

※ 2015년 8월 이후 지가변동률은 미고시 상태이며 2015년 7월 지가변동률과 동일하게 변동하는 것으로 추정함 \par
2) 오피스 자본수익률(단위: \%) \par
{\small\centering
\resizebox{\ifdim\width>\linewidth\linewidth\else\width\fi}{!}{%
\begin{tabular}{|c|c|c|}
\hline
\makecell{기간} & \makecell{2015년 1분기} & \makecell{2015년 2분기} \\
\hline
\makecell{자본수익률} & \makecell{0.35} & \makecell{0.30} \\
\hline
\end{tabular}}\par}

※ 2015년 2분기 이후 오피스 자본수익률은 미고시 상태이며 2015년 2분기 오피스 자본수익률과 동일하게 변동하는 것으로 추정함 \par
4. 거래사례 및 임대사례는 모두 정상적인 사례로 판단되며, 토지거래사례는 그 밖의 요인 보정치 산정에 활용함 \par
5. 건물의 감가수정은 정액법을 적용하며, 가치형성요인과 관련된 보정치는 계량분석내용을 준용하여 적용 가능함 \par
6. 유의수준이란 가설검증시 제1종 오류를 범할 확률의 허용한계, 즉, 오차 가능성을 말함 \par
\end{longquestion}
